%%%%%%%%%%%%%%%%%%%%%%%%%%%%%%%%%%%%%%%%%
% fphw Assignment
% LaTeX Template
% Version 1.0 (27/04/2019)
%
% This template originates from:
% https://www.LaTeXTemplates.com
%
% Authors:
% Class by Felipe Portales-Oliva (f.portales.oliva@gmail.com) with template 
% content and modifications by Vel (vel@LaTeXTemplates.com)
%
% Template (this file) License:
% CC BY-NC-SA 3.0 (http://creativecommons.org/licenses/by-nc-sa/3.0/)
%
%%%%%%%%%%%%%%%%%%%%%%%%%%%%%%%%%%%%%%%%%

%----------------------------------------------------------------------------------------
%	PACKAGES AND OTHER DOCUMENT CONFIGURATIONS
%----------------------------------------------------------------------------------------

\documentclass[
  12pt,
  %fleqn, % Default font size, values between 10pt-12pt are allowed
  %letterpaper, % Uncomment for US letter paper size
  %spanish, % Uncomment for Spanish
]{fphw}

% Template-specific packages
\usepackage[utf8]{inputenc} % Required for inputting international characters
\usepackage[T1]{fontenc} % Output font encoding for international characters
\usepackage{mathpazo} % Use the Palatino font

\usepackage{graphicx} % Required for including images

\usepackage{booktabs} % Required for better horizontal rules in tables

\usepackage{listings} % Required for insertion of code

\usepackage{enumerate} % To modify the enumerate environment

\usepackage{enumitem}

\usepackage{amsmath}

\usepackage{mathtools,amsthm}

\usepackage{float}

\usepackage{listings}

\usepackage{hyperref}

\lstdefinelanguage{R}{
  keywords={if, else, repeat, while, function, for, in, next, break},
  otherkeywords={TRUE, FALSE, NULL, NA, Inf, NaN},
  sensitive=true,
  morecomment=[l]\#,
  morestring=[b]",
}

\newcommand{\E}{\mathbb{E}}
\newcommand{\Var}{\operatorname{Var}}
\newcommand{\R}{\mathbb{R}}
\newcommand{\norm}[1]{\left\lVert #1 \right\rVert}


%----------------------------------------------------------------------------------------
%	ASSIGNMENT INFORMATION
%----------------------------------------------------------------------------------------

\title{Homework \#5} % Assignment title

\author{Shizhe Zhang} % Student name

\date{\today} % Due date

\institute{University of California, Berkeley \\ Department of Statistics} % Institute or school name

\class{EECS 227AT} % Course or class name

\professor{Gireeja Ranade} % Professor or teacher in charge of the assignment

%----------------------------------------------------------------------------------------

\begin{document}

\maketitle % Output the assignment title, created automatically using the information in the custom commands above

%----------------------------------------------------------------------------------------
%	ASSIGNMENT CONTENT
%----------------------------------------------------------------------------------------
\section*{Problem 1. Properties of the Frobenius Norm}
The Frobenius norm of a matrix $P$ is defined as
\[
\|P\|_F = \sqrt{\langle P, P\rangle} =
\sqrt{\sum_{1\le i\le m}\sum_{1\le j\le n} P_{ij}^2}, \tag{1}
\]
where for two matrices $P, Q \in \mathbb{R}^{m\times n}$, the canonical inner product defined over this space is
$\langle P, Q\rangle := \operatorname{tr}\!\left(P^\top Q\right) = \sum_{i,j} P_{ij}Q_{ij}$.
The previous definition of the inner product is equivalent to interpreting the matrices
$P$ and $Q$ as vectors of length $mn$ and taking the vector inner product of the respective $mn$-dimensional vectors.
The Cauchy-Schwarz inequality for the inner product follows in a straightforward way from the Cauchy-Schwarz
inequality for vectors:
\[
\langle P, Q\rangle =
\sum_{1\le i\le m}\sum_{1\le j\le n} P_{ij}Q_{ij}
\le
\left(\sum_{1\le i\le m}\sum_{1\le j\le n} P_{ij}^2\right)^{1/2}
\left(\sum_{1\le i\le m}\sum_{1\le j\le n} Q_{ij}^2\right)^{1/2}
= \|P\|_F \, \|Q\|_F. \tag{2}
\]

\begin{enumerate}
\item[(a)] Show that the Frobenius norm satisfies all three properties of a norm.

HINT: The easiest way to do this problem is to not look at the individual components $P_{ij}$, but instead use
the inner product formulation $\|P\|_F = \sqrt{\langle P, P\rangle}$.

\item[(b)] Write the Frobenius norm squared in terms of singular values.

HINT: The cyclic property of traces might be helpful: $\operatorname{tr}(PQR) = \operatorname{tr}(RPQ) = \operatorname{tr}(QRP)$.

\item[(c)] Express the Frobenius norm squared in terms of the $\ell_2$-norm of the columns of $P$ with $\vec p_i$ denoting column
$i$. Concretely, prove $\|P\|_F^2 = \sum_{i=1}^n \|\vec p_i\|_2^2$ where $\vec p_i$ are the columns of $P$.
\end{enumerate}

\subsection*{Answer}
\paragraph{(a)} % (fold)
1. Non-negativity: $\|P\|_F \geq 0$ for all $P$, and $\|P\|_F = 0$ if and only  if $P = 0$.

2. Absolute scalability: $\|\alpha P\|_F = |\alpha| \|P\|_F$ for all scalars $\alpha$ and matrices $P$.

It is easy to verify the first two properties.

3. Triangle inequality: $\|P + Q\|_F \leq \|P\|_F + \|Q\|_F$ for all matrices $P$ and $Q$.  

\begin{align*}
\|P + Q\|_F^2 &= \langle P + Q, P + Q\rangle \\
&= \langle P, P\rangle + 2\langle P, Q\rangle + \langle Q, Q\rangle \\
&\leq \|P\|_F^2 + 2\|P\|_ F \|Q\|_F + \|Q\|_F^2 = (\|P\|_F + \|Q\|_F)^2
\end{align*}

% paragraph (a) (end)

\paragraph{(b)} % (fold)
Let $P = U\Sigma V^\top$ be an SVD of $P$. Then
\begin{align*}
\|P\|_F^2 &= \langle P, P\rangle = \operatorname{tr}(P^\top P) = \operatorname{tr}(V\Sigma^\top U^\top U\Sigma V^\top) \\
&= \operatorname{tr}(V\Sigma^\top \Sigma V^\top) = \operatorname{tr}(\Sigma^\top \Sigma) = \sum_{i=1}^r \sigma_i^2,
\end{align*}
where $r$ is the rank of $P$ and $\sigma_i$ are the singular values of $P$. 

% paragraph (b) (end)

\paragraph{(c)} % (fold)
Let $\vec p_i$ denote the $i$-th column of $P$. Then
\begin{align*}
\|P\|_F^2 &= \langle P, P\rangle = \operatorname{tr}(P^\top P) = \operatorname{tr}\!\left(\begin{bmatrix}
\vec p_1^\top\\
\vec p_2^\top\\
\vdots\\
\vec p_n^\top
\end{bmatrix}\begin{bmatrix}
\vec p_1 & \vec p_2 & \cdots & \vec p_n
\end{bmatrix}\right) \\
&= \operatorname{tr}\!\left(\begin{bmatrix}
\vec p_1^\top \vec p_1 & \vec p_1^\top \vec p_2 & \cdots & \vec p_1^\top \vec p_n\\
\vec p_2^\top \vec p_1 & \vec p_2^\top \vec p_2 & \cdots & \vec p_2^\top \vec p_n\\
\vdots & \vdots & \ddots & \vdots\\
\vec p_n^\top \vec p_1 & \vec p_n^\top \vec p_2 & \cdots & \vec p_n^\top \vec p_n
\end{bmatrix}\right) = \sum_{i=1}^n \|\vec p_i\|_2^2.
\end{align*}


% paragraph (c) (end)



%----------------------------------------------------------------------------------------
\newpage
\section*{Problem 2. FTLA, SVD, Pseudoinverse, and Least-Squares}

Let $A \in \mathbb{R}^{m\times n}$ be a matrix, and let $\vec y \in \mathbb{R}^m$.
Let $r \doteq \operatorname{rank}(A)$, and let
\[
A = U\Sigma V^\top
=
\begin{bmatrix}
U_r & U_{m-r}
\end{bmatrix}
\begin{bmatrix}
\Sigma_r & 0_{r\times (n-r)}\\
0_{(m-r)\times r} & 0_{(m-r)\times (n-r)}
\end{bmatrix}
\begin{bmatrix}
V_r^\top\\
V_{n-r}^\top
\end{bmatrix}
=
U_r \Sigma_r V_r^\top
=
\sum_{i=1}^r \sigma_i \vec u_i \vec v_i^\top
\tag{3}
\]
be an SVD of $A$. In this problem, we will prove some properties about the SVD, and then apply them to derive
a unique solution to an optimization problem which generalizes minimum-norm and least-squares.

\begin{enumerate}
\item[(a)] Prove that $\mathcal{R}(A) = \mathcal{R}(U_r)$ and that $\mathcal{R}\!\left(A^\top\right) = \mathcal{R}(V_r)$.

\item[(b)] Use the fundamental theorem of linear algebra to prove that $\mathcal{N}(A) = \mathcal{R}(V_{n-r})$ and $\mathcal{N}\!\left(A^\top\right) = \mathcal{R}(U_{m-r})$.
\end{enumerate}

Now, we will use the SVD to derive the unique solution to the least-norm least-squares problem. That is, we
want to use the SVD to find the unique vector in $\mathbb{R}^n$ which solves the least-squares problem with minimum norm.
More formally, we wish to solve the least-norm least-squares problem:
\[
\min_{\vec x\in S}\ \|\vec x\|_2^2
\quad\text{where}\quad
S \doteq \arg\min_{\vec x\in \mathbb{R}^n}\ \|A\vec x - \vec y\|_2^2.
\tag{4}
\]
In words, we want to find the minimum-norm vector $\vec x$ in the set of all least-squares solutions $S$.%
\footnote{Here, the argmin is a set -- that is, it is the set of minimizers of the least-squares objective. A simpler example of this notation is that if
$f(x)\doteq 0$ for all $x \in \mathbb{R}$, then $\arg\min_{x\in\mathbb{R}} f(x) = \mathbb{R}$. More information about this notation is contained in section 1.3 of the course reader.}
This problem generalizes both the traditional least-squares and minimum-norm problems.

To solve the least-norm least-squares problem, we will use the SVD to define a concept called the pseudoinverse.
The matrix $A^\dagger \in \mathbb{R}^{n\times m}$ is called a pseudoinverse (sometimes a Moore-Penrose pseudoinverse) of $A$, where
\[
A^\dagger = V\Sigma^\dagger U^\top
=
\begin{bmatrix}
V_r & V_{n-r}
\end{bmatrix}
\begin{bmatrix}
\Sigma_r^{-1} & 0_{r\times (m-r)}\\
0_{(n-r)\times r} & 0_{(n-r)\times (m-r)}
\end{bmatrix}
\begin{bmatrix}
U_r^\top\\
U_{m-r}^\top
\end{bmatrix}
=
V_r \Sigma_r^{-1} U_r^\top
=
\sum_{i=1}^r \frac{1}{\sigma_i}\,\vec v_i \vec u_i^\top.
\tag{5}
\]
In this problem, we will show that $A^\dagger \vec y$ is the unique solution to the least-norm least-squares problem (4).

NOTE: In this problem, we do not assume that $A$ has full column rank or full row rank -- in fact, $A$ could even
be the matrix of all zeros -- so the least-squares solution we learn in class does not apply.

\begin{enumerate}
\item[(c)] Show that $\vec x \in S$ if and only if $A^\top A \vec x = A^\top \vec y$ (these are the so-called normal equations, which you may
remember from our discussion on least-squares).

\item[(d)] Show that $S = \{A^\dagger \vec y + \vec z \mid \vec z \in \mathcal{N}(A)\}$.

HINT: First, show that $A^\dagger \vec y \in S$. Then show that for any two $\vec x_1, \vec x_2 \in S$, that $\vec x_1 - \vec x_2 \in \mathcal{N}(A)$. Argue
that this implies the conclusion.

\item[(e)] Now show that $\{A^\dagger \vec y\} = \arg\min_{\vec x\in S} \|\vec x\|_2^2$.

HINT: Pick any $\vec z \in \mathcal{N}(A)$. Show that $A^\dagger \vec y$ is orthogonal to $\vec z$. Then show that
\[
\|A^\dagger \vec y\|_2^2 \le \|A^\dagger \vec y + \vec z\|_2^2,
\]
with equality if and only if $\vec z = \vec 0$. Argue that this implies the conclusion.
\end{enumerate}

\subsection*{Answer}
\begin{enumerate}
\item[(a)] Using the reduced SVD $A=U_r\Sigma_r V_r^\top$ with $\Sigma_r\in\mathbb{R}^{r\times r}$ invertible.

\emph{First, $\mathcal{R}(A)\subseteq \mathcal{R}(U_r)$.}
For any $\vec x\in\mathbb{R}^n$,
\[
A\vec x = U_r\Sigma_r V_r^\top \vec x \in \mathcal{R}(U_r),
\]
so $\mathcal{R}(A)\subseteq \mathcal{R}(U_r)$.

\emph{Second, $\mathcal{R}(U_r)\subseteq \mathcal{R}(A)$.}
Take any $\vec w\in\mathcal{R}(U_r)$. Then $\vec w=U_r\vec a$ for some $\vec a\in\mathbb{R}^r$.
Let $\vec b \doteq \Sigma_r^{-1}\vec a$ and define $\vec x \doteq V_r \vec b\in\mathbb{R}^n$.
Then $V_r^\top \vec x = V_r^\top V_r \vec b = \vec b$, hence
\[
A\vec x = U_r\Sigma_r V_r^\top \vec x = U_r\Sigma_r \vec b = U_r\vec a = \vec w.
\]
Thus $\vec w\in\mathcal{R}(A)$, proving $\mathcal{R}(U_r)\subseteq \mathcal{R}(A)$.
Therefore $\mathcal{R}(A)=\mathcal{R}(U_r)$.

For $A^\top$, note $A^\top = V_r\Sigma_r U_r^\top$. The same argument (swap $U_r$ and $V_r$) gives
$\mathcal{R}(A^\top)=\mathcal{R}(V_r)$.

\item[(b)] Since $V$ is orthogonal, its columns form an orthonormal basis of $\mathbb{R}^n$, and
\[
\mathbb{R}^n=\mathcal{R}(V_r)\oplus \mathcal{R}(V_{n-r}),\qquad
\mathbb{R}^m=\mathcal{R}(U_r)\oplus \mathcal{R}(U_{m-r}).
\]
From (a), $\mathcal{R}(A^\top)=\mathcal{R}(V_r)$. By the fundamental theorem of linear algebra,
\[
\mathcal{N}(A)=\mathcal{R}(A^\top)^\perp = \mathcal{R}(V_r)^\perp = \mathcal{R}(V_{n-r}).
\]
Similarly, from (a), $\mathcal{R}(A)=\mathcal{R}(U_r)$, and FTLA gives
\[
\mathcal{N}(A^\top)=\mathcal{R}(A)^\perp = \mathcal{R}(U_r)^\perp = \mathcal{R}(U_{m-r}).
\]

\item[(c)] Let $f(\vec x)\doteq \|A\vec x-\vec y\|_2^2=(A\vec x-\vec y)^\top(A\vec x-\vec y)$.
This is a convex quadratic, and $\vec x\in S$ iff $\nabla f(\vec x)=\vec 0$.
Compute
\[
\nabla f(\vec x)=2A^\top(A\vec x-\vec y).
\]
Thus $\nabla f(\vec x)=0$ iff $A^\top(A\vec x-\vec y)=0$, i.e.
\[
A^\top A\vec x = A^\top \vec y.
\]
Therefore $\vec x\in S \iff A^\top A\vec x=A^\top \vec y$.

\item[(d)]

By (c), it suffices to show $A^\top A(A^\dagger \vec y)=A^\top \vec y$.
Using $A=U_r\Sigma_r V_r^\top$ and $A^\dagger=V_r\Sigma_r^{-1}U_r^\top$,
\[
A^\top A A^\dagger
= (V_r\Sigma_r U_r^\top)(U_r\Sigma_r V_r^\top)(V_r\Sigma_r^{-1}U_r^\top)
= V_r\Sigma_r (U_r^\top U_r)\Sigma_r (V_r^\top V_r)\Sigma_r^{-1}U_r^\top
= A^\top.
\]
Multiplying by $\vec y$ gives $A^\top A(A^\dagger\vec y)=A^\top \vec y$, hence $A^\dagger \vec y\in S$.

If $\vec x_1, \vec x_2 \in S$, then $A^\top A\vec x_1=A^\top \vec y=A^\top A\vec x_2$, so $A^\top A(\vec x_1-\vec x_2)=0$.
Since $\mathcal{N}(A^\top A)=\mathcal{N}(A)$, $\vec x_1-\vec x_2 \in \mathcal{N}(A)$.

Thus any $\vec x\in S$ has the form $\vec x=A^\dagger \vec y + \vec z$ for some $\vec z\in\mathcal{N}(A)$, so
\[S \subseteq \{A^\dagger \vec y + \vec z \mid \vec z\in\mathcal{N}(A)\}.\] 

On the other hand, for any $\vec z\in\mathcal{N}(A)$, $A^\top A(A^\dagger \vec y + \vec z) = A^\top A A^\dagger \vec y + A^\top A \vec z = A^\top \vec y + 0 = A^\top \vec y$, so $A^\dagger \vec y + \vec z \in S$.
Thus $\{A^\dagger \vec y + \vec z \mid \vec z\in\mathcal{N}(A)\} \subseteq S$.

Combining gives $S = \{A^\dagger \vec y + \vec z \mid \vec z\in\mathcal{N}(A)\}$.


\item[(e)] By (d), any $\vec x\in S$ has the form $\vec x=A^\dagger\vec y+\vec z$ with $\vec z\in\mathcal{N}(A)$.
We show $A^\dagger\vec y \perp \mathcal{N}(A)$.

Let $\vec z\in\mathcal{N}(A)=\mathcal{R}(V_{n-r})$ (by (b)), so $\vec z=V_{n-r}\vec t$ for some $\vec t$ and hence
$V_r^\top \vec z = 0$ (orthogonality of $V_r$ and $V_{n-r}$).
Also $A^\dagger\vec y=V_r\Sigma_r^{-1}U_r^\top \vec y \in \mathcal{R}(V_r)$.
Therefore
\[
\langle A^\dagger\vec y,\vec z\rangle
= (V_r\Sigma_r^{-1}U_r^\top \vec y)^\top \vec z
= ( \Sigma_r^{-1}U_r^\top \vec y)^\top (V_r^\top \vec z)
= 0.
\]
Hence for $\vec x=A^\dagger\vec y+\vec z$,
\[
\|\vec x\|_2^2
= \|A^\dagger\vec y+\vec z\|_2^2
= \|A^\dagger\vec y\|_2^2 + \|\vec z\|_2^2
\ge \|A^\dagger\vec y\|_2^2,
\]
with equality iff $\|\vec z\|_2^2=0$, i.e. $\vec z=\vec 0$.
Thus the unique minimizer of $\|\vec x\|_2^2$ over $S$ is $A^\dagger\vec y$, i.e.
\[
\{A^\dagger\vec y\}=\arg\min_{\vec x\in S}\|\vec x\|_2^2.
\]
\end{enumerate}


%----------------------------------------------------------------------------------------
\newpage
\section*{Problem 3. Singular Values and Condition Number}

\begin{enumerate}
\item[(a)] Let $A \in \mathbb{R}^{m\times n}$ be a matrix. Show that
\[
\max_{\vec v:\ \|\vec v\|_2\le 1} \|A\vec v\|_2 = \sigma_{\max}\{A\}, \tag{6}
\]
where $\sigma_{\max}\{A\}$ is the largest singular value of $A$.

\item[(b)] Let $A \in \mathbb{R}^{n\times n}$ be an invertible matrix, let $\vec y \in \mathbb{R}^n$ be a vector, and let $\Delta \vec y \in \mathbb{R}^n$ be a vector (conceptually,
we think of it as having small norm). Let $\vec x$ and $\Delta \vec x$ be such that
\[
A\vec x = \vec y \tag{7}
\]
\[
A(\vec x + \Delta \vec x) = \vec y + \Delta \vec y. \tag{8}
\]
Show that
\[
\frac{\|\Delta \vec x\|_2}{\|\vec x\|_2}
\le \kappa(A)\cdot \frac{\|\Delta \vec y\|_2}{\|\vec y\|_2}, \tag{9}
\]
where $\kappa(A) \doteq \frac{\sigma_{\max}\{A\}}{\sigma_{\min}\{A\}}$ is called the condition number of the matrix $A$. It reflects how much the estimate
can change with respect to the perturbations $\Delta \vec y$ in the measurement $\vec y$.
\end{enumerate}


\subsection*{Answer}

\begin{enumerate}
\item[(a)] Let $A = U\Sigma V^\top$ be an SVD of $A$, where $\Sigma\in\mathbb{R}^{m\times n}$ has diagonal entries
$\sigma_1\ge \sigma_2\ge \cdots \ge \sigma_r > 0$ (and zeros thereafter), and $U,V$ are orthogonal (with suitable sizes).
For any $\vec v$ with $\|\vec v\|_2\le 1$, write $\vec v = V\vec z$ where $\vec z = V^\top \vec v$.
Since $V$ is orthogonal, $\|\vec z\|_2 = \|\vec v\|_2 \le 1$. Then
\[
\|A\vec v\|_2
= \|U\Sigma V^\top \vec v\|_2
= \|U\Sigma \vec z\|_2
= \|\Sigma \vec z\|_2,
\]
using that $U$ is orthogonal and thus preserves the $\ell_2$ norm.
Now $\Sigma \vec z$ has coordinates $(\sigma_1 z_1,\ldots,\sigma_r z_r, 0,\ldots)$, so
\[
\|\Sigma \vec z\|_2^2
= \sum_{i=1}^r \sigma_i^2 z_i^2
\le \sigma_1^2 \sum_{i=1}^r z_i^2
\le \sigma_1^2 \sum_{i=1}^n z_i^2
= \sigma_1^2 \|\vec z\|_2^2
\le \sigma_1^2.
\]
Hence $\|A\vec v\|_2 \le \sigma_1$ for all $\|\vec v\|_2\le 1$, so
\[
\max_{\|\vec v\|_2\le 1}\|A\vec v\|_2 \le \sigma_1 = \sigma_{\max}\{A\}.
\]
For achievability, take $\vec v = \vec v_1$, the first right singular vector (the first column of $V$), which has
$\|\vec v_1\|_2=1$. Then
\[
A\vec v_1 = U\Sigma V^\top \vec v_1 = U\Sigma \vec e_1 = U(\sigma_1 \vec e_1) = \sigma_1 \vec u_1,
\]
so $\|A\vec v_1\|_2 = \sigma_1$. Therefore the maximum equals $\sigma_1=\sigma_{\max}\{A\}$, proving (6).

\item[(b)] From (7)--(8),
\[
A(\vec x+\Delta \vec x) - A\vec x = (\vec y+\Delta \vec y) - \vec y
\quad\Longrightarrow\quad
A\,\Delta \vec x = \Delta \vec y.
\]
Since $A$ is invertible, $\Delta \vec x = A^{-1}\Delta \vec y$ and also $\vec x = A^{-1}\vec y$.
Using the operator norm induced by $\|\cdot\|_2$,
\[
\|\Delta \vec x\|_2 = \|A^{-1}\Delta \vec y\|_2 \le \|A^{-1}\|_2 \,\|\Delta \vec y\|_2,
\qquad
\|\vec y\|_2 = \|A\vec x\|_2 \le \|A\|_2\,\|\vec x\|_2.
\]
Rearranging the second inequality gives $\|\vec x\|_2 \ge \|\vec y\|_2/\|A\|_2$. Combining,
\[
\frac{\|\Delta \vec x\|_2}{\|\vec x\|_2}
\le
\frac{\|A^{-1}\|_2\,\|\Delta \vec y\|_2}{\|\vec y\|_2/\|A\|_2}
=
\|A\|_2\,\|A^{-1}\|_2 \cdot \frac{\|\Delta \vec y\|_2}{\|\vec y\|_2}.
\]
By part (a) applied to $A$ and to $A^{-1}$, $\|A\|_2=\sigma_{\max}(A)$ and $\|A^{-1}\|_2=\sigma_{\max}(A^{-1})$.
If $A$ is invertible with singular values $\sigma_1\ge\cdots\ge\sigma_n>0$, then $A^{-1}$ has singular values
$1/\sigma_1\le\cdots\le 1/\sigma_n$, hence $\sigma_{\max}(A^{-1}) = 1/\sigma_{\min}(A)$.
Therefore
\[
\|A\|_2\,\|A^{-1}\|_2
=
\sigma_{\max}(A)\cdot \frac{1}{\sigma_{\min}(A)}
=
\frac{\sigma_{\max}(A)}{\sigma_{\min}(A)}
=
\kappa(A),
\]
and substituting yields (9):
\[
\frac{\|\Delta \vec x\|_2}{\|\vec x\|_2}
\le \kappa(A)\cdot \frac{\|\Delta \vec y\|_2}{\|\vec y\|_2}.
\]
\end{enumerate}


%----------------------------------------------------------------------------------------
\newpage
\section*{Problem 4. PCA and low-rank compression}

We have a data matrix
\[
X =
\begin{bmatrix}
\vec x_1^\top\\
\vec x_2^\top\\
\vdots\\
\vec x_n^\top
\end{bmatrix}
\]
of size $n\times d$ containing $n$ data points\footnote{Data matrices are sometimes represented as above, and sometimes as the transpose of the matrix here. Make sure you always check this, and
recall that based on the definition of the data matrix, the definition of the covariance matrix also changes.}
$\vec x_1, \vec x_2, \ldots, \vec x_n$, with $\vec x_i \in \mathbb{R}^d$. Note
that $\vec x_i^\top$ is the $i$th row of $X$. Assume that the data matrix is centered, i.e. each column of $X$ is zero mean. In this
problem, we will show equivalence between the following three problems:

\textbf{(P1)} Finding a line going through the origin that maximizes the variance of the scalar projections of the points
on the line. Formally P1 solves the problem:
\[
\arg\max_{\vec u\in\mathbb{R}^d:\ \vec u^\top \vec u = 1}\ \vec u^\top C \vec u \tag{10}
\]
with $C = \frac{1}{n}\sum_{i=1}^n \vec x_i \vec x_i^\top$ denoting the covariance matrix associated with the centered data.

\textbf{(P2)} Finding a line going through the origin that minimizes the sum of squares of the $\ell_2$ distances from the
points to their vector projections. Formally P2 solves the minimization problem:
\[
\arg\min_{\vec u\in\mathbb{R}^d:\ \vec u^\top \vec u = 1}\ \sum_{i=1}^n \min_{v_i\in\mathbb{R}} \|\vec x_i - v_i \vec u\|_2^2. \tag{11}
\]
Note that the vector projection of $\vec x$ on $\vec u$ is given by $v^\star \vec u$, where
\[
v^\star = \arg\min_{v\in\mathbb{R}} \|\vec x - v\vec u\|_2^2, \tag{12}
\]
and we will show that $v^\star = \langle \vec x, \vec u\rangle$ in part (a).

\textbf{(P3)} Finding a rank-one approximation to the data matrix. Formally P3 solves the minimization problem:
\[
\arg\min_{Y:\ \operatorname{rank}(Y)\le 1}\ \|X - Y\|_F. \tag{13}
\]

Note that loosely speaking, two problems are said to be ``equivalent" if the solution of one can be ``easily"
translated to the solution of the other. Some form of ``easy" translations include adding/subtracting a constant or
some quantity depending on the data points.

Note the significance of these results. P1 is finding the first principal component of $X$, the direction that
maximizes variance of scalar projections. P2 says that this direction also minimizes the distances between the
points to their vector projections along this direction. If we view the distances as errors in approximating the
points by their projections along a line, then the error is minimized by choosing the line in the same direction
as the first principal component. Finally P3 tells us that finding a rank one matrix to best approximate the data
matrix (in terms of error computed using Frobenius norm) is equivalent to finding the first principal component
as well!

\begin{enumerate}
\item[(a)] Consider the line $L = \{\vec x_0 + a\vec u : a \in \mathbb{R}\}$, with $\vec x_0 \in \mathbb{R}^d$, $\vec u^\top \vec u = 1$. Recall that the vector projection of a
point $\vec x \in \mathbb{R}^d$ on to the line $L$ is given by $\vec z = \vec x_0 + a^\star \vec u$, where $a^\star$
is given by:
\[
a^\star = \arg\min_a \|\vec x_0 + a\vec u - \vec x\|_2. \tag{14}
\]
Show that $a^\star = (\vec x - \vec x_0)^\top \vec u$. Use this to show that the square of the distance between $x$ and its vector
projection on $L$ is given by:
\[
\|\vec x - \vec z\|_2^2 = \|\vec x - \vec x_0\|_2^2 - \big((\vec x - \vec x_0)^\top \vec u\big)^2. \tag{15}
\]

\item[(b)] Show that P2 is equivalent to P1.

HINT: Start with P2 and using the result from part (a) show that it is equivalent to P1.

\item[(c)] Show that every matrix $Y \in \mathbb{R}^{n\times d}$ with rank at most $1$, can be expressed as $Y = \vec v \vec u^\top$ for some $\vec v \in \mathbb{R}^n$,
$\vec u \in \mathbb{R}^d$ and $\|\vec u\|_2 = 1$.

\item[(d)] Show that P3 is equivalent to P2.

HINT: Use the result from part (c) to show that P3 is equivalent to:
\[
\arg\min_{\vec u\in\mathbb{R}^d:\ \vec u^\top\vec u = 1,\ \vec v\in\mathbb{R}^n}\ \|X - \vec v \vec u^\top\|_F^2. \tag{16}
\]
Prove that this is equivalent to P2.
\end{enumerate}


\subsection*{Answer}
\begin{enumerate}
\item[(a)] Define
\[
f(a) \doteq \|\vec x_0 + a\vec u - \vec x\|_2^2.
\]
Since $\vec u^\top \vec u = 1$, expand:
\begin{align*}
f(a)
&= (\vec x_0 + a\vec u - \vec x)^\top(\vec x_0 + a\vec u - \vec x) \\
&= \|\vec x_0-\vec x\|_2^2 + 2a\,\vec u^\top(\vec x_0-\vec x) + a^2\|\vec u\|_2^2 \\
&= \|\vec x-\vec x_0\|_2^2 - 2a\,(\vec x-\vec x_0)^\top \vec u + a^2.
\end{align*}
Differentiate and set to zero:
\[
f'(a) = -2(\vec x-\vec x_0)^\top \vec u + 2a = 0
\quad\Longrightarrow\quad
a^\star = (\vec x-\vec x_0)^\top \vec u.
\]
Let $\vec z = \vec x_0 + a^\star \vec u$. Plugging $a^\star$ into the expansion gives the minimum value:
\[
\|\vec x-\vec z\|_2^2
= f(a^\star)
= \|\vec x-\vec x_0\|_2^2 - \big((\vec x-\vec x_0)^\top \vec u\big)^2,
\]
which is (15).

\item[(b)] In P2, for each $i$ and fixed $\vec u$ with $\|\vec u\|_2=1$, the inner minimization
\[
\min_{v_i\in\mathbb{R}} \|\vec x_i - v_i \vec u\|_2^2
\]
is the squared distance from $\vec x_i$ to the line $\{v\vec u: v\in\mathbb{R}\}$ (i.e., $L$ with $\vec x_0=\vec 0$).
Applying (15) with $\vec x_0=\vec 0$ and $\vec x=\vec x_i$ yields
\[
\min_{v_i\in\mathbb{R}} \|\vec x_i - v_i \vec u\|_2^2
= \|\vec x_i\|_2^2 - (\vec x_i^\top \vec u)^2.
\]
Therefore the P2 objective is equivalent to
\[
\sum_{i=1}^n \min_{v_i} \|\vec x_i - v_i \vec u\|_2^2
= \sum_{i=1}^n \|\vec x_i\|_2^2 - \sum_{i=1}^n (\vec x_i^\top \vec u)^2.
\]
Eliminating the constant $\sum_{i=1}^n \|\vec x_i\|_2^2$ (which does not depend on $\vec u$) from the objective, P2 is equivalent to
\[
\arg\min_{\|\vec u\|_2=1}\ \sum_{i=1}^n \min_{v_i} \|\vec x_i - v_i \vec u\|_2^2
=
\arg\max_{\|\vec u\|_2=1}\ \sum_{i=1}^n (\vec x_i^\top \vec u)^2.
\]
Now compute
\[
\sum_{i=1}^n (\vec x_i^\top \vec u)^2
= \sum_{i=1}^n \vec u^\top (\vec x_i \vec x_i^\top)\vec u
= \vec u^\top \left(\sum_{i=1}^n \vec x_i \vec x_i^\top\right)\vec u
= n\,\vec u^\top C \vec u,
\]
where $C=\frac1n\sum_{i=1}^n \vec x_i\vec x_i^\top$. Hence P2 is equivalent to P1.

\item[(c)] Let $Y\in\mathbb{R}^{n\times d}$ with $\operatorname{rank}(Y)\le 1$.

i) $\operatorname{rank}(Y)=0$, take $\vec v=\vec 0$ and any unit $\vec u$.

ii) $\operatorname{rank}(Y)=1$, so the column space of $Y$ is one-dimensional:
there exists $\vec v\neq \vec 0$ such that every column of $Y$ is a scalar multiple of $\vec v$.
\begin{align*}
\Longrightarrow Y 
& = \begin{bmatrix} \alpha_1 \vec v & \alpha_2 \vec v & \cdots & \alpha_d \vec v \end{bmatrix} \\
& = \vec v \begin{bmatrix} \alpha_1 & \alpha_2 & \cdots & \alpha_d \end{bmatrix} \\
& = \vec v\,\vec \alpha^\top \\
& = (\|\vec \alpha\|_2\vec v)\left(\frac{\vec \alpha}{\|\vec \alpha\|_2}\right)^\top \\
& = \tilde{\vec v}\,\vec u^\top,
\end{align*}

\item[(d)] Using (c), the constraint $\operatorname{rank}(Y)\le 1$ is equivalent to $Y=\vec v\vec u^\top$ for some
$\vec v\in\mathbb{R}^n$, $\vec u\in\mathbb{R}^d$ with $\|\vec u\|_2=1$. Thus P3 is equivalent to
\[
\arg\min_{\|\vec u\|_2=1,\ \vec v\in\mathbb{R}^n}\ \|X-\vec v\vec u^\top\|_F,
\]
Which is equivalent to maximizing the square of the objective, so P3 is equivalent to
\[
\|X-\vec v\vec u^\top\|_F^2
= \sum_{i=1}^n \|\vec x_i^\top - v_i \vec u^\top\|_2^2
= \sum_{i=1}^n \|\vec x_i - v_i \vec u\|_2^2.
\]
\[
\min_{\vec v\in\mathbb{R}^n}\ \sum_{i=1}^n \|\vec x_i - v_i \vec u\|_2^2
= \sum_{i=1}^n \min_{v_i\in\mathbb{R}} \|\vec x_i - v_i \vec u\|_2^2.
\]
Thus
\[
\arg\min_{\|\vec u\|_2=1,\ \vec v\in\mathbb{R}^n}\ \|X-\vec v\vec u^\top\|_F^2
=
\arg\min_{\|\vec u\|_2=1}\ \sum_{i=1}^n \min_{v_i\in\mathbb{R}} \|\vec x_i - v_i \vec u\|_2^2,
\]
which is exactly P2. Hence P3 is equivalent to P2.
\end{enumerate}



%----------------------------------------------------------------------------------------
\newpage
\section*{Problem 5. Operator Norms}
For a matrix $A \in \mathbb{R}^{m\times n}$, the induced norm or operator norm $\|A\|_p$
is defined as
\[
\|A\|_p \doteq \max_{\vec x\ne \vec 0}\ \frac{\|A\vec x\|_p}{\|\vec x\|_p}. \tag{17}
\]
In this problem, we provide a characterization of the induced norm for certain values of $p$.
Let $a_{ij}$ denote the $(i, j)$-th entry of $A$. Prove the following:

\begin{enumerate}
\item[(a)] $\|A\|_2 = \sigma_{\max}\{A\}$, the maximum singular value of $A$.

HINT: Consider connecting $\|A\|_2^2$ to a particular Rayleigh coefficient.

\item[(b)] $\|A\|_1$ is the maximum absolute column sum of $A$,
\[
\|A\|_1 = \max_{1\le j\le n}\ \sum_{i=1}^m |a_{ij}|. \tag{18}
\]
HINT: Write $A\vec x$ as a linear combination of the columns of $A$ to obtain $\|A\vec x\|_1 = \left\|\sum_{i=1}^n x_i\cdot \vec a_i\right\|_1$,
where $\vec a_i$ denotes the $i$-th column of $A$. Then apply triangle inequality to terms within the sum.

\item[(c)] $\|A\|_\infty$ is the maximum absolute row sum of $A$,
\[
\|A\|_\infty = \max_{1\le i\le m}\ \sum_{j=1}^n |a_{ij}|. \tag{19}
\]
HINT: First write $\|A\vec x\|_\infty = \max_i \left|\sum_{j=1}^n a_{ij}x_j\right|$. Then apply triangle inequality and use the fact that
$|x_j| \le \max_i |x_i|,\ \forall j$.
\end{enumerate}


\subsection*{Answer}

\begin{enumerate}
\item[(a)] By definition,
\[
\|A\|_2 \doteq \max_{\vec x\neq 0}\frac{\|A\vec x\|_2}{\|\vec x\|_2}
= \max_{\|\vec x\|_2=1}\|A\vec x\|_2.
\]
Hence
\[
\|A\|_2^2
= \max_{\|\vec x\|_2=1}\|A\vec x\|_2^2
= \max_{\|\vec x\|_2=1} (A\vec x)^\top (A\vec x)
= \max_{\|\vec x\|_2=1} \vec x^\top (A^\top A)\vec x.
\]
The quantity $\vec x^\top (A^\top A)\vec x$ over $\|\vec x\|_2=1$ is the Rayleigh quotient of the symmetric
positive semidefinite matrix $A^\top A$, whose maximum equals its largest eigenvalue:
\[
\max_{\|\vec x\|_2=1} \vec x^\top (A^\top A)\vec x = \lambda_{\max}(A^\top A).
\]
If $\sigma_1\ge \sigma_2\ge \cdots \ge 0$ are the singular values of $A$, then the eigenvalues of $A^\top A$
are $\sigma_i^2$, and thus $\lambda_{\max}(A^\top A)=\sigma_{\max}(A)^2$. Therefore
\[
\|A\|_2^2 = \sigma_{\max}(A)^2
\quad\Longrightarrow\quad
\|A\|_2 = \sigma_{\max}(A).
\]

\item[(b)] Let $\vec a_j\in\mathbb{R}^m$ denote the $j$th column of $A$. For any $\vec x\in\mathbb{R}^n$,
\[
A\vec x = \sum_{j=1}^n x_j \vec a_j.
\]
Using the triangle inequality for $\|\cdot\|_1$,
\[
\|A\vec x\|_1
= \left\|\sum_{j=1}^n x_j \vec a_j\right\|_1
\le \sum_{j=1}^n |x_j|\,\|\vec a_j\|_1.
\]
Let $M \doteq \max_{1\le j\le n}\|\vec a_j\|_1$ (the maximum absolute column sum). Then
\[
\|A\vec x\|_1 \le M\sum_{j=1}^n |x_j| = M\|\vec x\|_1,
\]
so for all $\vec x\neq 0$,
\[
\frac{\|A\vec x\|_1}{\|\vec x\|_1} \le M
\quad\Longrightarrow\quad
\|A\|_1 \le M.
\]
To show equality, choose an index $j^\star \in \arg\max_j \|\vec a_j\|_1$ and set $\vec x = \vec e_{j^\star}$
(the $j^\star$th standard basis vector). Then $\|\vec x\|_1=1$ and $A\vec x = \vec a_{j^\star}$, hence
\[
\frac{\|A\vec x\|_1}{\|\vec x\|_1} = \|\vec a_{j^\star}\|_1 = M.
\]
Thus $\|A\|_1 \ge M$, and combined with $\|A\|_1\le M$ we get
\[
\|A\|_1 = \max_{1\le j\le n}\|\vec a_j\|_1 = \max_{1\le j\le n}\sum_{i=1}^m |a_{ij}|.
\]

\item[(c)] For any $\vec x\in\mathbb{R}^n$,
\[
\|A\vec x\|_\infty
= \max_{1\le i\le m}\left|\sum_{j=1}^n a_{ij}x_j\right|.
\]
Fix $i$. By triangle inequality and the bound $|x_j|\le \|\vec x\|_\infty$ for all $j$,
\[
\left|\sum_{j=1}^n a_{ij}x_j\right|
\le \sum_{j=1}^n |a_{ij}|\,|x_j|
\le \left(\sum_{j=1}^n |a_{ij}|\right)\|\vec x\|_\infty.
\]
Thus,
\[
\|A\vec x\|_\infty
\le \left(\max_{1\le i\le m}\sum_{j=1}^n |a_{ij}|\right)\|\vec x\|_\infty.
\]
Let $R \doteq \max_{1\le i\le m}\sum_{j=1}^n |a_{ij}|$ (maximum absolute row sum). Then for all $\vec x\neq 0$,
\[
\frac{\|A\vec x\|_\infty}{\|\vec x\|_\infty} \le R
\quad\Longrightarrow\quad
\|A\|_\infty \le R.
\]
For the reverse inequality, pick $i^\star\in\arg\max_i \sum_{j=1}^n |a_{ij}|$ and define $\vec x\in\mathbb{R}^n$ by
\[
x_j \doteq \operatorname{sign}(a_{i^\star j}) \in \{-1,0,1\},
\]
with the convention $\operatorname{sign}(0)=0$. Then $\|\vec x\|_\infty = 1$, and the $i^\star$th component satisfies
\[
(A\vec x)_{i^\star} = \sum_{j=1}^n a_{i^\star j} \operatorname{sign}(a_{i^\star j})
= \sum_{j=1}^n |a_{i^\star j}|.
\]
Hence
\[
\|A\vec x\|_\infty \ge |(A\vec x)_{i^\star}| = \sum_{j=1}^n |a_{i^\star j}| = R,
\]
so
\[
\|A\|_\infty \ge \frac{\|A\vec x\|_\infty}{\|\vec x\|_\infty} \ge R.
\]
Combining with $\|A\|_\infty \le R$ gives
\[
\|A\|_\infty = \max_{1\le i\le m}\sum_{j=1}^n |a_{ij}|.
\]
\end{enumerate}

%----------------------------------------------------------------------------------------
\newpage
\section*{Problem 6. Homework Process}

With whom did you work on this homework? List the names and SIDs of your group members.

NOTE: If you didn’t work with anyone, you can put “none” as your answer.


\subsection*{Answer}
none

%----------------------------------------------------------------------------------------

\end{document}
